% ======================================================================
% capitulos/linea3/e3_sgsi.tex
% ----------------------------------------------------------------------
% LÍNEA CE03 — INFRAESTRUCTURA TECNOLÓGICA
% Semestre 2 · Entregable 1 de 3: PLANIFICACIÓN E IMPLEMENTACIÓN DEL SGSI
% Códigos: CE0321, CE0322, CE0323, CE0324.
% Estructura según la plantilla oficial S2_E1_Planificacion_e_Implementacion_del_SGSI.
% ALCANCE: infraestructura REAL en operación (servidor dockerizado de la oficina,
% publicación de wispfibra.com por Cloudflare Tunnel + DNS, VPS con el hub
% WireGuard y 13 zonas).
% Las tablas de activos, riesgos, tratamiento, SoA y trazabilidad se GENERAN con
% sgsi/sgsi_datos.py (misma fuente que el Excel sgsi/WISP_SGSI_Matrices.xlsx).
% Los documentos de soporte del equipo están en los Anexos A–F (anexos_sgsi.tex).
% ======================================================================
\chapter{Planificación e Implementación del SGSI}
\label{chap:ce03-sgsi}

\entregable{CE03 --- Infraestructura Tecnológica · competencia CE032 Gestión de la Seguridad de la Información}
{Semestre 2 · Entregable 1: Planificación e Implementación del SGSI}
{CE0321 Planificación de seguridad · CE0322 Implementación de controles · CE0323 Monitoreo y mejora · CE0324 Ética ACM}
{Completo. Pendiente: firmas de la declaración ética y del acta de revisión por la dirección.}

\begin{table}[htbp]
\centering
\caption{Productos obligatorios del entregable y dónde se encuentran.}
\label{tab:ce03-sgsi-productos}
\small
\begin{tabularx}{\textwidth}{L{4.4cm}XL{2.2cm}}
\toprule
\textbf{Producto (rúbrica)} & \textbf{Ubicación} & \textbf{Formato} \\
\midrule
Informe del SGSI & Este capítulo & PDF \\
Inventario de activos & \cref{tab:ce03-sgsi-activos}; hoja \emph{Inventario} de \texttt{sgsi/WISP\_SGSI\_Matrices.xlsx} & XLSX + PDF \\
Matriz de evaluación y tratamiento de riesgos & \cref{tab:ce03-sgsi-riesgos,tab:ce03-sgsi-tratamiento}; hojas \emph{Riesgos} y \emph{Tratamiento} & XLSX + PDF \\
Declaración de Aplicabilidad (SoA) & \cref{tab:ce03-sgsi-soa}; hoja \emph{SoA} & XLSX + PDF \\
Políticas y procedimientos del SGSI & \cref{sec:ce03-sgsi-politicas} y Anexos A, C, D y E & PDF \\
Declaración de ética de seguridad & \cref{sec:ce03-sgsi-etica} y Anexo F (código de ética) & PDF firmado \\
\bottomrule
\end{tabularx}
\end{table}

% ======================================================================
\section{Resumen ejecutivo del SGSI}

\begin{itemize}
  \item \textbf{Alcance:} la infraestructura que WISP opera hoy: un servidor dockerizado en
  la oficina central que aloja el sistema de gestión y su base de datos, publicado hacia
  internet mediante un túnel de salida y registros DNS de Cloudflare (sin puertos de
  entrada en el router de la oficina); la VPS que
  concentra la VPN WireGuard; y los 13 routers y OLT de las zonas, cuya segmentación real (VLAN 10 de transporte GPON/PPPoE)
  se tomó de la configuración de SUCASCO. Protege los datos de
  unos 455 clientes, los cobros, las credenciales y la configuración de red.
  \item \textbf{Riesgos principales:} 20 riesgos sobre 26 activos (ISO/IEC~27005,
  probabilidad $\times$ impacto). Siete son \textbf{altos}: routers de zona sin filtrado
  (20, verificado en la configuración real de SUCASCO), phishing (16), incidentes no
  detectados (16), fuga de datos personales (15), pérdida de datos por respaldos solo
  locales (15), robo de secretos (15) y falla del servidor único (15).
  \item \textbf{Controles:} la Declaración de Aplicabilidad tiene 37 entradas del Anexo~A de
  ISO/IEC~27001:2022: 5 implementadas, 20 parciales, 10 planificadas y 2 exclusiones
  justificadas.
  \item \textbf{PHVA:} \emph{Planificar} (contexto, alcance, riesgos, SoA);
  \emph{Hacer} (políticas, procedimiento de accesos, planes de incidentes, continuidad y
  concientización, y controles de base de datos ya probados); \emph{Verificar}
  (auditoría interna del 18/11/2026: 7 no conformidades, 2 observaciones y 5
  conformidades); \emph{Actuar} (acciones correctivas y revisión por la dirección).
  \item \textbf{Conclusión:} el sistema ya tiene buenos controles de aplicación ---acceso
  por rol, secretos fuera del repositorio, servicios internos no expuestos y publicación
  sin puertos abiertos---. La mayor exposición está fuera del software: routers de zona
  sin filtrado, respaldos que no salen del servidor y un único equipo que concentra todo.
\end{itemize}

% ======================================================================
\section{Contexto de la organización}

WISP Ingeniería en Redes y Telecomunicaciones es una MYPE de Juliaca que brinda internet
por fibra óptica (FTTH/GPON) y televisión digital en 13 zonas rurales de Puno (Cabana,
Cabanilla, Cabanillas, Capachica, Choquehuanca, Domingo, Huata, Mañazo, Muñani, Ñuñoa,
Pucará, Sucasco y Vilque), con unos 455 clientes y 11 colaboradores. La regulan OSIPTEL y
el MTC y, como responsable de bancos de datos personales, la Ley N.\textsuperscript{o}~29733.
El análisis estratégico completo está en el \cref{chap:diagnostico}.

\paragraph{Qué protege el SGSI.} La confidencialidad de los datos personales y de cobro;
la integridad del padrón, de los cobros y de la configuración de red, porque un dato
erróneo corta a un cliente al día; y la disponibilidad de la gestión y del servicio de
internet de las 13 zonas.

\paragraph{Cuestiones internas.} Una sola persona concentra la administración; la
madurez de seguridad es inicial; todo el sistema corre en un único equipo de escritorio.
\paragraph{Cuestiones externas.} Dependencia de un único proveedor mayorista y de
proveedores de nube (Cloudflare, VPS); cortes de energía y de fibra en zonas rurales;
obligación de notificar incidentes de datos personales y de seguridad digital; ataques
conocidos contra RouterOS expuesto a internet.

\subsection{Infraestructura dentro del alcance}

\begin{figure}[H]
\centering
\begin{tikzpicture}[x=1cm,y=1cm,scale=0.92,every node/.style={transform shape}]
  \node[cajaext,text width=2.6cm] (usr) at (-6,2.2) {Personal, clientes y visitantes (HTTPS)};
  \node[cajaext,text width=2.6cm] (cf) at (-2.4,2.2) {Cloudflare: DNS de wispfibra.com + túnel};
  \node[caja,text width=4.4cm,fill=upeublue!10] (srv) at (-2.4,-0.6) {\textbf{Servidor de oficina (Docker)}\\cloudflared · nginx · backend\\frontend · web · whatsapp\\postgres (127.0.0.1)};
  \node[caja,text width=2.8cm] (vl) at (-6.6,-0.6) {Router de la oficina (sin puertos de entrada)};
  \node[caja,text width=2.7cm] (vps) at (2.8,2.2) {VPS: hub WireGuard\\wg0 · UDP 51820};
  \node[caja,text width=2.7cm] (zn) at (2.8,-0.6) {13 routers MikroTik + OLT; VLAN 10 GPON/PPPoE};
  \node[cajaext,text width=2.6cm] (ext) at (2.8,-2.8) {WhatsApp, APIs Perú};
  \draw[flecha] (usr) -- (cf);
  \draw[flecha] (srv.north) -- node[font=\tiny,right,align=left]{conexión\\de salida} (cf.south);
  \draw[upeublue,thick] (vl) -- (srv);
  \draw[tunel] (srv.east) -- node[font=\tiny,above,sloped]{WireGuard} (vps.south west);
  \draw[tunel] (vps) -- node[font=\tiny,right]{WireGuard / SSTP} (zn);
  \draw[flecha] (srv.south east) -- (ext.west);
\end{tikzpicture}
\caption{Infraestructura real dentro del alcance del SGSI.}
\label{fig:ce03-sgsi-infra}
\end{figure}

\begin{table}[htbp]
\centering
\caption{Exposición a internet: qué se publica y qué queda cerrado.}
\label{tab:ce03-sgsi-publicacion}
\small
\begin{tabularx}{\textwidth}{L{3.6cm}L{3.4cm}X}
\toprule
\textbf{Servicio} & \textbf{Cómo se expone} & \textbf{Control} \\
\midrule
Sitio, panel y API & \texttt{https://wispfibra.com} (y \texttt{www}): \texttt{/} sitio, \texttt{/app} panel, \texttt{/api} API (la usa la app móvil) $\rightarrow$ túnel $\rightarrow$ \texttt{nginx:80} & HTTPS lo termina Cloudflare; el router de la oficina no abre ningún puerto de entrada; Nginx agrega HSTS y cabeceras de seguridad. \\
PostgreSQL, WhatsApp, panel interno & No se publican & Escuchan solo en 127.0.0.1 del servidor. \\
API en el puerto 8000 & No se publica & El backend usa la red del host: el firewall del servidor debe cerrar el 8000 salvo desde la red de Docker (R-13). \\
Hub VPN & IP pública de la VPS, UDP 51820 & Solo peers con llave registrada; administración de la VPS por SSH con llave. \\
Routers de zona & WAN del ISP local & Hoy sin filtrado (R-01); objetivo: gestión solo por WireGuard. \\
Gestión de la OLT & 192.168.100.0/24 en \texttt{ether9} & Alcanzable desde la LAN de clientes por falta de filtro forward (R-20). \\
\bottomrule
\end{tabularx}
\end{table}

\noindent El diseño híbrido del semestre~1 (host VMware ESXi y VPC de AWS,
\cref{chap:ce03-red,chap:ce03-dc}) es la evolución planificada de esta infraestructura; sus
componentes se incorporarán al inventario y a la SoA cuando se implementen.

\subsection{Partes interesadas y sus expectativas de seguridad}

\begin{table}[H]
\centering
\caption{Partes interesadas del SGSI.}
\label{tab:ce03-sgsi-interesados}
\small
\begin{tabularx}{\textwidth}{L{3.2cm}X}
\toprule
\textbf{Parte interesada} & \textbf{Expectativa de seguridad} \\
\midrule
Clientes (unos 455) & Sus datos (DNI, ubicación, fotos de vivienda, pagos) no se filtran; no los cortan por error; se les avisa si hay una brecha. \\
Gerencia & Continuidad del negocio, cumplimiento legal y control de quién hace qué. \\
Administrador de sistemas y red & Acceso seguro y trazable al servidor, a la VPS y a los 13 routers; procedimientos escritos. \\
Técnicos, vendedores y atención & Herramientas que protejan su cuenta; reglas claras sobre qué datos pueden ver y guardar. \\
Proveedor mayorista & Que la red de WISP no origine tráfico malicioso. \\
Cloudflare, proveedor de la VPS, Meta & Uso conforme a sus términos; responsabilidad compartida sobre la configuración. \\
ANPD (MINJUSDH) & Bancos de datos inscritos, medidas de seguridad y notificación de incidentes de datos personales. \\
CNSD (PCM) y OSIPTEL & Reporte de incidentes críticos de seguridad digital y registro de interrupciones. \\
Docente y universidad & Pruebas éticas y sin datos personales reales en los entregables. \\
\bottomrule
\end{tabularx}
\end{table}

% ======================================================================
\section{Alcance del SGSI}

\subsection{Declaración de alcance}
El SGSI comprende la infraestructura tecnológica de WISP en operación: el servidor de la
oficina central con sus siete contenedores Docker (\texttt{wisp\_backend},
\texttt{wisp\_postgres}, \texttt{wisp\_frontend}, \texttt{wisp\_web},
\texttt{wisp\_nginx}, \texttt{wisp\_whatsapp} y \texttt{wisp\_cloudflared}), sus volúmenes de
datos, evidencias y respaldos; la publicación de \texttt{wispfibra.com} mediante el túnel y el DNS de Cloudflare; la VPS con el hub WireGuard; los 13 routers MikroTik con su VLAN 10 de transporte y la red de gestión de la OLT, y
las 13 OLT; la app móvil; y los procesos de venta, instalación, cobranza, cortes, soporte y
administración de la red que dependen de ellos.

\subsection{Límites y exclusiones}
\begin{itemize}
  \item \textbf{Límites físicos:} oficina central de Juliaca (servidor, router, switch,
  UPS, estaciones) y gabinetes de las 13 zonas.
  \item \textbf{Límites tecnológicos:} los 26 activos del inventario
  (\cref{tab:ce03-sgsi-activos}).
  \item \textbf{Límites organizacionales:} todo el personal de WISP y los terceros con
  acceso lógico.
  \item \textbf{Exclusión --- red óptica pasiva (fibra, mufas, cajas NAP):} no procesa
  información; su protección física la cubre la gestión de planta.
  \item \textbf{Exclusión --- equipos del cliente (ONU, WiFi doméstico):} están bajo control
  del cliente.
  \item \textbf{Exclusión --- infraestructura de Cloudflare, del proveedor de la VPS y de
  Meta:} la operan terceros; dentro del alcance queda la \emph{configuración} y las cuentas
  de WISP en esos servicios.
  \item \textbf{Exclusión --- host ESXi y recursos de AWS del diseño del semestre~1:} aún no
  existen; entrarán al alcance cuando se implementen.
\end{itemize}

% ======================================================================
\section{Liderazgo y objetivos de seguridad}

\subsection{Política de seguridad de la información}
\label{sec:ce03-sgsi-objetivos-pol}

\begin{upeubox}{Política general del SGSI de WISP (PO-SGSI-01, versión 1)}
La Gerencia de WISP se compromete a proteger la confidencialidad, integridad y
disponibilidad de la información de sus clientes y de la operación de su red. Para ello:
(1) el acceso a la información y a la red se concede por rol y con el mínimo privilegio;
(2) ningún servicio interno se expone a internet: la publicación es solo por el túnel y la
administración de la red solo por la VPN; (3) la información crítica se respalda fuera del
servidor y su restauración se prueba; (4) los incidentes se reportan en un máximo de 30
minutos y se notifican a las autoridades en los plazos de ley; (5) los datos personales se
tratan conforme a la Ley N.\textsuperscript{o}~29733; y (6) el SGSI se revisa al menos una vez
al año y después de cada incidente grave. Esta política es obligatoria para todo el
personal y los terceros con acceso.
\end{upeubox}

\subsection{Roles y responsabilidades}

Se separan el rol que administra los sistemas y el que los supervisa (segregación de
funciones, A.5.3), tal como propuso el plan de incidentes del equipo:

\begin{table}[htbp]
\centering
\caption{Roles del SGSI.}
\label{tab:ce03-sgsi-roles}
\small
\begin{tabularx}{\textwidth}{L{3.6cm}L{3cm}X}
\toprule
\textbf{Rol} & \textbf{Asignado a} & \textbf{Responsabilidad} \\
\midrule
Alta dirección & Gerencia & Aprueba la política y la SoA, asigna recursos y preside el Comité de Seguridad y la revisión por la dirección. \\
Oficial de seguridad de la información & Rol a tiempo parcial designado por la Gerencia & Coordina el SGSI, mantiene riesgos y SoA, lidera el CSIRT, audita accesos y notifica a reguladores. \\
Responsable de infraestructura & Administrador de sistemas y red & Opera el servidor, la red de oficina, la VPS, el túnel y los routers. \\
Responsable de sistemas & Equipo de desarrollo / administrador & Opera el sistema, la base de datos y sus respaldos. \\
Suplente técnico & Técnico capacitado con cuenta de emergencia & Cubre al administrador ante ausencia (R-04). \\
Usuarios & Técnicos, vendedores, atención & Cumplen la política de uso aceptable y reportan incidentes. \\
\bottomrule
\end{tabularx}
\end{table}

\subsection{Objetivos de seguridad medibles}

Se mantienen los objetivos OS-1 a OS-6 del SGSI de la Unidad~I y se agregan tres derivados
de la infraestructura real:

\begin{table}[htbp]
\centering
\caption{Objetivos de seguridad, indicador, meta y responsable.}
\label{tab:ce03-sgsi-objetivos}
\small
\begin{tabularx}{\textwidth}{L{0.9cm}XL{3.3cm}L{1.9cm}L{2cm}}
\toprule
\textbf{ID} & \textbf{Objetivo} & \textbf{Indicador} & \textbf{Meta} & \textbf{Responsable} \\
\midrule
OS-1 & Controlar el acceso según el rol & Cuentas nominales con rol asignado & 100\,\% & Oficial \\
OS-2 & Cifrar las comunicaciones & Zonas gestionadas por WireGuard (no SSTP) & 13 de 13 & Infraestructura \\
OS-3 & Garantizar la recuperación & Restauraciones exitosas en las pruebas trimestrales & 100\,\%; RTO BD $\leq$ 2 h & Sistemas \\
OS-4 & Trazabilidad de operaciones críticas & Altas, cortes y pagos con registro de autor & 100\,\% & Sistemas \\
OS-5 & Reducir los riesgos altos & Riesgos altos con residual medio o bajo & 7 de 7 & Oficial \\
OS-6 & Cultura de seguridad & Personal capacitado (Anexo E) & 11 de 11 & Oficial \\
OS-7 & Proteger accesos administrativos & Cuentas con MFA: Cloudflare, dominio, VPS y panel admin & 100\,\% & Infraestructura \\
OS-8 & Cero exposición no autorizada & Puertos de entrada abiertos fuera de la lista aprobada (\cref{tab:ce03-sgsi-publicacion}) & 0 & Infraestructura \\
OS-9 & Respaldos fuera del servidor & Días con copia externa verificada / días del mes & 100\,\% & Sistemas \\
\bottomrule
\end{tabularx}
\end{table}

% ======================================================================
\section{Planificación}

\subsection{Inventario de activos y clasificación de la información}
\label{sec:ce03-sgsi-inventario}

Escala C-I-D: A (alta), M (media), B (baja). El mismo inventario está en la hoja
\emph{Inventario} del libro \texttt{sgsi/WISP\_SGSI\_Matrices.xlsx}.

{\small
\begin{longtable}{L{0.8cm}L{4.0cm}L{2.2cm}L{2.6cm}C{0.3cm}C{0.3cm}C{0.3cm}L{1.3cm}L{1.1cm}}
\caption{Inventario de activos del SGSI.}\label{tab:ce03-sgsi-activos}\\
\toprule
\textbf{ID} & \textbf{Activo} & \textbf{Propietario} & \textbf{Ubicación} & \textbf{C} & \textbf{I} & \textbf{D} & \textbf{Clasif.} & \textbf{Critic.} \\
\midrule
\endfirsthead
\toprule
\textbf{ID} & \textbf{Activo} & \textbf{Propietario} & \textbf{Ubicación} & \textbf{C} & \textbf{I} & \textbf{D} & \textbf{Clasif.} & \textbf{Critic.} \\
\midrule
\endhead
\bottomrule
\endlastfoot
\csname @@input\endcsname capitulos/linea3/tablas/activos.tex
\end{longtable}
}

\subsection{Evaluación de riesgos (ISO/IEC 27005)}
\label{sec:ce03-sgsi-riesgos}

\paragraph{Metodología.} Igual que en el SGSI de la Unidad~I: probabilidad (P) e impacto
(I) de 1 a 5; nivel = P $\times$ I; \textbf{bajo} 1--6, \textbf{medio} 8--12, \textbf{alto}
15--25. Se aceptan riesgos residuales bajos o medios con aprobación de la Gerencia. Las
vulnerabilidades marcadas con \textsuperscript{v} se \textbf{verificaron} en la
configuración real (router de SUCASCO, \texttt{docker-compose.yml}, código del backend y
repositorio). La línea base de la Unidad~I ---22 riesgos sobre los activos del sistema
heredado--- está en el Anexo~B.

{\small
\begin{longtable}{L{0.9cm}L{1.1cm}L{3.2cm}L{4.6cm}C{0.4cm}C{0.4cm}C{0.6cm}L{1.0cm}}
\caption{Evaluación de riesgos: activo, amenaza, vulnerabilidad, probabilidad, impacto y nivel.}\label{tab:ce03-sgsi-riesgos}\\
\toprule
\textbf{ID} & \textbf{Activo} & \textbf{Amenaza} & \textbf{Vulnerabilidad} & \textbf{P} & \textbf{I} & \textbf{P$\times$I} & \textbf{Nivel} \\
\midrule
\endfirsthead
\toprule
\textbf{ID} & \textbf{Activo} & \textbf{Amenaza} & \textbf{Vulnerabilidad} & \textbf{P} & \textbf{I} & \textbf{P$\times$I} & \textbf{Nivel} \\
\midrule
\endhead
\bottomrule
\endlastfoot
\csname @@input\endcsname capitulos/linea3/tablas/riesgos.tex
\end{longtable}
}

\noindent Resultado: 7 riesgos altos, 12 medios y 1 bajo.

\subsection{Tratamiento de riesgos}
\label{sec:ce03-sgsi-tratamiento}

{\small
\begin{longtable}{L{0.9cm}L{1.4cm}L{5.6cm}L{2.0cm}L{1.4cm}L{1.6cm}}
\caption{Plan de tratamiento: estrategia, acciones, responsable, plazo y riesgo residual.}\label{tab:ce03-sgsi-tratamiento}\\
\toprule
\textbf{ID} & \textbf{Estrategia} & \textbf{Acciones} & \textbf{Responsable} & \textbf{Plazo} & \textbf{Residual} \\
\midrule
\endfirsthead
\toprule
\textbf{ID} & \textbf{Estrategia} & \textbf{Acciones} & \textbf{Responsable} & \textbf{Plazo} & \textbf{Residual} \\
\midrule
\endhead
\bottomrule
\endlastfoot
\csname @@input\endcsname capitulos/linea3/tablas/tratamiento.tex
\end{longtable}
}

\noindent Con el plan aplicado ningún riesgo residual queda alto. Se aceptan formalmente
R-12 y R-14 (residual 9, medio) y R-17 (6, bajo).

\subsection{Declaración de Aplicabilidad (SoA)}
\label{sec:ce03-sgsi-soa}

\begin{figure}[H]
\centering
\begin{tikzpicture}[x=1cm,y=1cm]
  \node[cajasec,text width=2.8cm] (r) at (0,0) {\textbf{R-01}\\Acceso no autorizado a routers de zona};
  \node[caja,text width=2.6cm] (c) at (3.8,0) {\textbf{A.8.20 · A.8.22 · A.8.9}\\Redes, segregación, configuración};
  \node[caja,text width=3.0cm] (i) at (7.7,0) {\textbf{Implementación}\\Plantilla endurecida: descarte por defecto, gestión por WireGuard};
  \node[cajaext,text width=2.6cm] (e) at (11.5,0) {\textbf{Prueba / evidencia}\\PC-05, PC-06 · export del router};
  \draw[flecha] (r) -- (c); \draw[flecha] (c) -- (i); \draw[flecha] (i) -- (e);
\end{tikzpicture}
\caption{Trazabilidad riesgo $\rightarrow$ control $\rightarrow$ implementación $\rightarrow$ evidencia (riesgo mayor).}
\label{fig:ce03-sgsi-traza}
\end{figure}

Estados: \textbf{Impl.} funcionando hoy; \textbf{Parcial} existe pero no cubre todo el
alcance; \textbf{Plan.} se implementa durante el semestre.

{\small
\begin{longtable}{L{1.4cm}L{3.5cm}L{1.7cm}L{1.2cm}L{6.2cm}}
\caption{Declaración de Aplicabilidad (Anexo A de ISO/IEC 27001:2022).}\label{tab:ce03-sgsi-soa}\\
\toprule
\textbf{Control} & \textbf{Nombre} & \textbf{Riesgos} & \textbf{Estado} & \textbf{Justificación / evidencia} \\
\midrule
\endfirsthead
\toprule
\textbf{Control} & \textbf{Nombre} & \textbf{Riesgos} & \textbf{Estado} & \textbf{Justificación / evidencia} \\
\midrule
\endhead
\bottomrule
\endlastfoot
\csname @@input\endcsname capitulos/linea3/tablas/soa.tex
\end{longtable}
}

\noindent Resumen: 35 entradas aplicables (5 implementadas, 20 parciales y 10
planificadas) y 2 exclusiones justificadas.

% ======================================================================
\section{Soporte --- políticas y procedimientos del SGSI}
\label{sec:ce03-sgsi-politicas}

\begin{table}[H]
\centering
\caption{Políticas y procedimientos del SGSI, vinculados a componentes reales.}
\label{tab:ce03-sgsi-politicas-tab}
\small
\begin{tabularx}{\textwidth}{L{3cm}X}
\toprule
\textbf{Documento} & \textbf{Reglas principales} \\
\midrule
PO-SGSI-01 Seguridad de la información & Política general (\cref{sec:ce03-sgsi-objetivos-pol}); revisión anual y tras incidentes graves. \\
PO-SGSI-02 Control de accesos & Cuentas nominales; prohibidas las cuentas genéricas (\texttt{admin}); Cloudflare, dominio y VPS con MFA y correo de recuperación de la empresa; el servidor de oficina se administra solo desde la red interna y por SSH con llave; routers solo por WireGuard; revisión trimestral; baja inmediata al cese. Se ejecuta con el procedimiento \textbf{PR-SI-001} (Anexo~A). \\
PR-SGSI-03 Gestión de incidentes & Reporte en $\leq$ 30 min, severidad, CSIRT, contención, erradicación, recuperación y lecciones aprendidas (\cref{sec:ce03-sgsi-incidentes} y Anexo~C). \\
PO-SGSI-04 Proveedores y nube & Cloudflare, proveedor de la VPS, Meta, APIs Perú y mayorista con términos revisados; inventario de cuentas y de quién las administra; ningún dato personal en servicios no evaluados; revisión anual de cada proveedor. \\
PO-SGSI-05 Respaldo y recuperación & 3-2-1: \texttt{pg\_dump} diario en el servidor, copia externa diaria cifrada y copia en disco externo; configuración de routers exportada cada noche; \texttt{.env} y \texttt{docker-compose.yml} respaldados cifrados; prueba de restauración trimestral (Anexo~D). \\
PO-SGSI-06 Uso aceptable & El personal no guarda fotos ni datos de clientes en su celular ni los envía por WhatsApp personal; no comparte su cuenta; bloquea el teléfono de trabajo; reporta su pérdida de inmediato. \\
\bottomrule
\end{tabularx}
\end{table}

% ======================================================================
\section{Operación --- implementación de controles de gestión}

\subsection{Gestión de identidades y accesos (gobernanza)}

\begin{table}[H]
\centering
\caption{Matriz de accesos por rol (mínimo privilegio).}
\label{tab:ce03-sgsi-iam}
\footnotesize
\setlength{\tabcolsep}{3pt}
\begin{tabularx}{\textwidth}{L{2.7cm}*{6}{>{\raggedright\arraybackslash}X}}
\toprule
\textbf{Rol} & \textbf{Cloudflare / VPS / dominio} & \textbf{Servidor Docker} & \textbf{Routers de zona} & \textbf{Router de oficina} & \textbf{Sistema} & \textbf{Base de datos} \\
\midrule
Gerencia & Propietario (MFA, emergencia) & --- & --- & --- & Admin & Lectura (vistas) \\
Administrador de sistemas y red & Admin (MFA) & Admin (sudo, SSH con llave) & Admin & Admin & Admin & Dueño (migraciones) \\
Oficial de seguridad & Lectura / auditor & Lectura de logs & Lectura & Lectura & Lectura & Lectura (vistas) \\
Suplente técnico & --- & Operador (cuenta de emergencia) & Admin (cuenta de emergencia) & --- & Técnico & --- \\
Técnico & --- & --- & --- & --- & Técnico (sus zonas) & --- \\
Vendedor de campo & --- & --- & --- & --- & Vendedor & --- \\
Atención al cliente & --- & --- & --- & --- & Atención & --- \\
Backend (servicio) & --- & Contenedor & SSH limitado (usuario \texttt{api}) & --- & --- & Rol de aplicación \\
Respaldos (servicio) & Escritura en el destino externo & Lectura del volumen & --- & --- & --- & Rol de respaldo \\
\bottomrule
\end{tabularx}
\end{table}

\begin{itemize}
  \item \textbf{MFA:} obligatoria en Cloudflare, en el registrador del dominio y en el
  proveedor de la VPS, porque quien controla el DNS o el túnel puede suplantar el panel
  (R-05). El panel de administración se protegerá con Cloudflare Access (código de un solo
  uso) mientras el sistema no tenga MFA propio.
  \item \textbf{Credenciales:} ningún secreto en el repositorio (el \texttt{.env.example} no
  tiene valores reales y el backend se niega a arrancar en producción con claves de
  ejemplo); \texttt{.env} con permisos 600; llaves WireGuard generadas en cada equipo y PSK
  por peer; rotación anual de llaves, del token del túnel y de las contraseñas de servicio;
  cuenta de emergencia sellada y rotada tras cada uso. El procedimiento completo está en el
  Anexo~A.
\end{itemize}

\subsection{Gestión de incidentes, continuidad y recuperación (planificación)}
\label{sec:ce03-sgsi-incidentes}

El plan completo de gestión de incidentes está en el Anexo~C y el plan de continuidad del
negocio con su análisis de impacto (BIA) en el Anexo~D. Sus elementos clave:

\begin{itemize}
  \item \textbf{Severidad:} crítico (datos personales comprometidos, 3 o más zonas o el
  servidor central caídos, ransomware), alto, medio y bajo, con notificación interna en 15
  min, 30 min, 2 h y 24 h, respectivamente.
  \item \textbf{CSIRT:} Gerente (líder de crisis), Oficial (coordinador), Administrador
  (analista técnico), técnicos (campo) y comercial (comunicación), cada uno con suplente.
  \item \textbf{Contención en esta infraestructura:} bloquear la cuenta en el sistema, cortar
  la publicación desactivando la ruta del túnel en Cloudflare, revocar el peer WireGuard
  afectado, detener el contenedor comprometido o aislar el router de zona.
  \item \textbf{Notificación externa:} ANPD (datos personales) y CNSD (incidentes críticos de
  seguridad digital) en 48~h, OSIPTEL por interrupciones y DIVINDAT ante indicios de delito.
\end{itemize}

\begin{table}[htbp]
\centering
\caption{Servicios críticos: objetivos de recuperación actuales y meta.}
\label{tab:ce03-sgsi-rto}
\small
\begin{tabularx}{\textwidth}{L{3.2cm}L{2.5cm}L{2.3cm}X}
\toprule
\textbf{Servicio} & \textbf{Hoy (RTO / RPO)} & \textbf{Meta (RTO / RPO)} & \textbf{Cómo se alcanza la meta} \\
\midrule
Base de datos & $\approx$ 4 h / hasta 7 días si se pierde el disco & 2 h / 1 h & Copia externa diaria y archivado WAL; restauración probada (en laboratorio: 650\,000 filas en 4 s, \cref{lst:ce02-respaldo}). \\
Sistema (API, panel, app) & 4--8 h (reinstalar en otro equipo) & 2 h & Equipo de respaldo con Docker y el compose listos; procedimiento de migración documentado. \\
Publicación (túnel) & 30 min & 30 min & Crear un token nuevo y apuntar el mismo nombre DNS. \\
Hub VPN & Sin procedimiento & 1 h & Script de reconstrucción y copia cifrada de \texttt{wg0.conf}. \\
Configuración de routers & Manual & 1 h & Exportación nocturna automática. \\
Internet por zona & 4 h (lógica) / 24 h (física) & Igual & Revertir configuración; equipo de repuesto homogéneo. \\
Atención al cliente & 2 h & 2 h & Línea telefónica directa. \\
\bottomrule
\end{tabularx}
\end{table}

% ======================================================================
\section{Ética de la seguridad de la información}
\label{sec:ce03-sgsi-etica}

El informe ético-legal completo del equipo ---marco legal, 34 obligaciones verificadas, 12
riesgos ético-legales, ocho políticas y el Código de Ética y Conducta de WISP--- está en el
Anexo~F. Aquí se declara la ética del equipo como profesionales de seguridad.

\begin{upeubox}{Declaración de ética del equipo del proyecto}
\textbf{Código adoptado.} El Código de Ética y Conducta Profesional de la ACM (2018), en
especial los principios 1.2 Evitar daños, 1.6 Respetar la privacidad, 1.7 Respetar la
confidencialidad, 2.5 Evaluar los riesgos de forma exhaustiva y 2.9 Diseñar sistemas
robustos y usablemente seguros; y los cuatro cánones de (ISC)\textsuperscript{2}: proteger a
la sociedad, actuar con honestidad y legalidad, prestar un servicio diligente y desarrollar
la profesión.

\textbf{Alcance ético de las pruebas.} Solo se prueban sistemas de WISP o del laboratorio
con autorización escrita de la Gerencia. No se ejecutan ataques de denegación de servicio ni
pruebas contra Cloudflare, el proveedor de la VPS, el mayorista u otros terceros. No se
prueba sobre una zona en producción fuera de una ventana de mantenimiento acordada.

\textbf{Consentimiento y autorización.} Antes de cada prueba se firma una autorización con
sistemas, fechas, técnicas y responsable; la firman la Gerencia de WISP y el docente.

\textbf{Confidencialidad y protección de datos.} Los datos personales se tratan conforme a
la Ley N.\textsuperscript{o}~29733 y su reglamento; en los entregables solo hay datos
anonimizados.

\textbf{Divulgación responsable.} Una vulnerabilidad en un producto de terceros (RouterOS,
OLT, librerías) se reporta en privado al fabricante y a la Gerencia, no se publica antes de
90 días o de su corrección y nunca se explota.

\textbf{Conflictos de interés.} Cada integrante declara cualquier relación laboral,
comercial o familiar con WISP o sus proveedores; quien audita un área no la administra.

\medskip
\begin{tabularx}{\linewidth}{XX}
\rule{5.5cm}{0.4pt} & \rule{5.5cm}{0.4pt} \\
Cabana Sulca, Cristian & Fernández Sánchez, Jherson Jean Piero \\[10pt]
\rule{5.5cm}{0.4pt} & \rule{5.5cm}{0.4pt} \\
Mamani Mamani, Alessandro Pastor & Montes Mamani, Andrés Lino \\
\end{tabularx}
\end{upeubox}

\paragraph{Coherencia con el comportamiento del proyecto.}
\begin{itemize}
  \item La configuración real de SUCASCO contenía nombres, DNI y direcciones de clientes;
  se analizó en privado y en el documento solo aparecen valores de red.
  \item La base de datos se probó con datos ficticios, no con una copia de producción.
  \item Las IP públicas, llaves y tokens se reemplazaron por marcadores como
  \texttt{<EIP-HUB>}.
  \item La exposición de fotos de clientes en el repositorio se reportó como no conformidad
  en lugar de omitirse (NC-02).
\end{itemize}

% ======================================================================
\section{Evaluación del desempeño --- monitoreo, medición y auditoría interna}

\subsection{Métricas e indicadores del SGSI}

\begin{table}[H]
\centering
\caption{Avance de los objetivos de seguridad (línea base al 18/11/2026).}
\label{tab:ce03-sgsi-metricas}
\small
\begin{tabularx}{\textwidth}{L{0.9cm}XL{4.2cm}L{1.8cm}}
\toprule
\textbf{ID} & \textbf{Indicador} & \textbf{Línea base} & \textbf{Meta} \\
\midrule
OS-1 & Cuentas nominales con rol & Sistema: 100\,\% (4 roles); routers y servidor: sin medir & 100\,\% \\
OS-2 & Zonas gestionadas por WireGuard & SUCASCO aún en SSTP; resto por relevar & 13 de 13 \\
OS-3 & Restauraciones probadas & 1 en laboratorio (4 s, integridad verificada) & 1 por trimestre \\
OS-4 & Operaciones críticas trazadas & Historial del cliente; auditoría de pagos en laboratorio & 100\,\% \\
OS-5 & Riesgos altos con residual medio o bajo & 0 de 7 (plan definido) & 7 de 7 \\
OS-6 & Personal capacitado & 0 de 11 & 11 de 11 \\
OS-7 & Cuentas privilegiadas con MFA & Sin medir & 100\,\% \\
OS-8 & Puertos de entrada no autorizados & 0 en la oficina (publicación por túnel); 13 routers expuestos & 0 \\
OS-9 & Días con copia externa & 0 (copia semanal manual) & 100\,\% \\
\bottomrule
\end{tabularx}
\end{table}

\subsection{Auditoría interna}

\begin{figure}[H]
\centering
\begin{tikzpicture}[x=1cm,y=1cm]
  \node[caja,text width=2.2cm] (a) at (0,0) {Criterios\\ISO/IEC 27001};
  \node[caja,text width=2.2cm] (b) at (2.9,0) {Lista de verificación};
  \node[caja,text width=2.2cm] (c) at (5.8,0) {Evidencia recolectada};
  \node[caja,text width=2.2cm] (d) at (8.7,0) {Conforme / no conforme};
  \node[cajasec,text width=2.2cm] (e) at (11.6,0) {Acción correctiva};
  \draw[flecha] (a) -- (b); \draw[flecha] (b) -- (c); \draw[flecha] (c) -- (d); \draw[flecha] (d) -- (e);
  \draw[flecha] (e.south) -- ++(0,-0.7) -| (a.south);
\end{tikzpicture}
\caption{Ciclo de auditoría interna del SGSI.}
\label{fig:ce03-sgsi-auditoria}
\end{figure}

\begin{table}[htbp]
\centering
\caption{Plan de la auditoría interna.}
\label{tab:ce03-sgsi-plan-auditoria}
\small
\begin{tabularx}{\textwidth}{L{2.6cm}X}
\toprule
Fecha & 18/11/2026 \\
Tipo & Documental y técnica (sin pruebas intrusivas) \\
Criterios & Cláusulas 4 a 10 de ISO/IEC 27001:2022 y controles de la SoA \\
Alcance & Repositorio del sistema (rama \texttt{main}), \texttt{docker-compose.yml} y configuración de Nginx del servidor de oficina, configuración exportada del router de SUCASCO y documentos del SGSI \\
\bottomrule
\end{tabularx}
\end{table}

{\small
\begin{longtable}{L{1.0cm}L{1.0cm}L{6.1cm}L{2.0cm}L{3.1cm}}
\caption{Hallazgos de la auditoría interna y acciones correctivas.}\label{tab:ce03-sgsi-hallazgos}\\
\toprule
\textbf{ID} & \textbf{Control} & \textbf{Hallazgo y evidencia} & \textbf{Tipo} & \textbf{Acción correctiva} \\
\midrule
\endfirsthead
\toprule
\textbf{ID} & \textbf{Control} & \textbf{Hallazgo y evidencia} & \textbf{Tipo} & \textbf{Acción correctiva} \\
\midrule
\endhead
\bottomrule
\endlastfoot
NC-01 & 8.20 & El router de SUCASCO no tiene reglas de filtrado y expone servicios de gestión (export del 27/08/2026). & NC mayor & R-01 (Infraestructura, 3 meses). \\
NC-02 & 5.34 & Hay 8 fotos de evidencias de clientes versionadas en Git (\texttt{backend/media/evidencias}). & NC mayor & R-10 (Sistemas, 1 semana). \\
NC-03 & 8.13 & El respaldo diario queda en el mismo disco del servidor; la copia externa es manual y semanal; no hay registro de restauraciones en producción. & NC menor & R-03 (Sistemas, 1 mes). \\
NC-04 & 8.24 & Credenciales PPPoE y de routers guardadas en texto plano (\texttt{clientes.password\_pppoe}, \texttt{zonas.password\_mikrotik}). & NC menor & R-02 (Sistemas, 2 meses). \\
NC-05 & 8.20 & El cliente SSH acepta cualquier llave de host (\texttt{paramiko.AutoAddPolicy}). & NC menor & R-08 (Sistemas, 1 mes). \\
NC-06 & 8.22 & En SUCASCO la VLAN 10 está puenteada con \texttt{ether6} y \texttt{ether7} y no hay filtro forward: la LAN de clientes llega a la gestión de la OLT (192.168.100.0/24). & NC menor & R-20 (Infraestructura, 1 mes). \\
NC-07 & 8.8 & La imagen \texttt{cloudflare/cloudflared:latest} no fija versión y el backend corre en la red del host con NET\_ADMIN. & NC menor & R-13 (Infraestructura, 1 mes). \\
OB-01 & 5.34 & El README y \texttt{scripts/deploy.sh} publican la IP real de la VPS. & Observación & Reemplazar por marcadores. \\
OB-02 & 8.2 & La terminal RouterOS usa lista negra de comandos. & Observación & R-16: lista blanca. \\
C-01 & 5.15 & Cada endpoint exige JWT y rol (\texttt{require\_roles}). & Conforme & --- \\
C-02 & 5.17 & No hay secretos en el repositorio; \texttt{.env.example} sin valores; el backend no arranca en producción con claves de ejemplo. & Conforme & --- \\
C-03 & 8.5 & Freno a la fuerza bruta en el login (8 intentos en 5 min, bloqueo de 15 min). & Conforme & --- \\
C-04 & 8.20 & PostgreSQL, WhatsApp y el panel escuchan solo en 127.0.0.1; la publicación es por túnel de salida, sin puertos de entrada en la oficina. & Conforme & --- \\
C-05 & 8.32 & Cambios del esquema en 32 migraciones Alembic revisables. & Conforme & --- \\
\end{longtable}
}

\subsection{Revisión por la dirección}

\begin{table}[H]
\centering
\caption{Acta de revisión por la dirección (equipo actuando como dirección).}
\label{tab:ce03-sgsi-revision}
\small
\begin{tabularx}{\textwidth}{L{3.2cm}X}
\toprule
Entradas revisadas & Auditoría interna (7 NC, 2 observaciones, 5 conformidades), objetivos OS-1 a OS-9, riesgos residuales, resultados de la prueba de restauración en laboratorio y cambios del contexto (infraestructura dockerizada en producción publicada en \texttt{wispfibra.com}). \\
Conclusión & El SGSI es \textbf{adecuado en su diseño} pero \textbf{aún no eficaz} en la red de zonas y en los respaldos, donde están las dos NC de mayor impacto. \\
Decisiones & (1) Resolver NC-02 en una semana y NC-01 empezando por SUCASCO; (2) activar MFA en Cloudflare, dominio y VPS de inmediato; (3) automatizar la copia externa diaria; (4) designar al Oficial de seguridad separado del administrador; (5) aceptar los residuales de R-12, R-14 y R-17; (6) presupuestar el servidor definitivo con disco espejo. \\
Próximo ciclo & Primera restauración en producción, ejercicio de mesa del playbook de fuga de datos (Anexo~C), primera campaña de concientización (Anexo~E) y nueva auditoría tras migrar las zonas. \\
Firmas & \rule{5cm}{0.4pt} (Gerencia) \qquad \rule{5cm}{0.4pt} (Oficial) \\
\bottomrule
\end{tabularx}
\end{table}

% ======================================================================
\section{Mejora continua --- matriz maestra de trazabilidad}
\label{sec:ce03-sgsi-matriz}

Las pruebas PC-xx son las de la matriz de conectividad del diseño de red
(\cref{tab:ce03-pruebas}); P6, P7 y B1--B4 son las pruebas de base de datos del entregable
de plataforma de datos (\cref{chap:ce02-datos}).

{\small
\begin{longtable}{L{1.0cm}L{1.9cm}L{4.4cm}L{2.8cm}L{3.4cm}}
\caption{Matriz maestra: riesgo $\rightarrow$ control $\rightarrow$ implementación $\rightarrow$ prueba $\rightarrow$ evidencia.}\label{tab:ce03-sgsi-maestra}\\
\toprule
\textbf{Riesgo} & \textbf{Control} & \textbf{Implementación} & \textbf{Prueba} & \textbf{Evidencia} \\
\midrule
\endfirsthead
\toprule
\textbf{Riesgo} & \textbf{Control} & \textbf{Implementación} & \textbf{Prueba} & \textbf{Evidencia} \\
\midrule
\endhead
\bottomrule
\endlastfoot
\csname @@input\endcsname capitulos/linea3/tablas/maestra.tex
\end{longtable}
}

% ======================================================================
\section{Resultados y conclusiones}

\begin{itemize}
  \item \textbf{Riesgos identificados y tratados:} 20 riesgos sobre 26 activos (7 altos, 12
  medios y 1 bajo); los 20 tienen estrategia, responsable y plazo, y ninguno queda alto
  después del tratamiento.
  \item \textbf{Riesgos residuales aceptados:} R-12 (cortes de fibra y energía), R-14
  (bloqueo de WhatsApp) y R-17 (robo de equipos en campo).
  \item \textbf{Controles funcionando:} acceso por rol, secretos fuera del repositorio,
  freno a la fuerza bruta, servicios internos no expuestos, publicación por túnel sin
  puertos de entrada y cambios versionados (C-01 a C-05). Probados en laboratorio: roles de
  base de datos, auditoría de pagos, triggers de integridad y respaldo con restauración.
  \item \textbf{Controles pendientes:} filtrado de los routers de zona, MFA, copia externa
  automática, centralización de logs, segregación de funciones y servidor con redundancia.
  \item \textbf{Auditoría interna:} 2 NC mayores, 5 NC menores, 2 observaciones y 5
  conformidades, con acción correctiva, responsable y plazo.
  \item \textbf{Limitaciones:} toda la operación depende hoy de un solo equipo de escritorio;
  la línea base de 12 de las 13 zonas requiere relevar su configuración.
  \item \textbf{Mejoras para el siguiente ciclo:} MFA propio en el panel, alerta automática
  por exportaciones masivas, inscripción de los bancos de datos ante la ANPD, migración al
  diseño híbrido del semestre~1 y evaluación del alineamiento formal con ISO/IEC~27001.
\end{itemize}

% ======================================================================
\section*{Referencias}
\addcontentsline{toc}{section}{Referencias}
{\small
\begin{enumerate}[label={[\arabic*]},leftmargin=2em]
  \item ISO/IEC 27001:2022, \emph{Information security management systems --- Requirements}. Ginebra: ISO, 2022.
  \item ISO/IEC 27002:2022, \emph{Information security controls}. Ginebra: ISO, 2022.
  \item ISO/IEC 27005:2022, \emph{Guidance on managing information security risks}. Ginebra: ISO, 2022.
  \item ISO/IEC 27035-1:2023, \emph{Information security incident management --- Principles and process}. Ginebra: ISO, 2023.
  \item ISO 22301:2019, \emph{Business continuity management systems --- Requirements}. Ginebra: ISO, 2019.
  \item P. Cichonski \emph{et al.}, ``Computer Security Incident Handling Guide'', NIST SP 800-61 Rev. 2, 2012.
  \item NIST, ``Building an Information Technology Security Awareness and Training Program'', NIST SP 800-50, 2003.
  \item INCIBE, ``Plan de Contingencia y Continuidad de Negocio'', guía para empresas.
  \item Congreso de la República del Perú, Ley N.\textsuperscript{o} 29733, Ley de Protección de Datos Personales, y su Reglamento, D.S. N.\textsuperscript{o} 016-2024-JUS.
  \item Congreso de la República del Perú, Ley N.\textsuperscript{o} 30096, Ley de Delitos Informáticos.
  \item ACM, ``ACM Code of Ethics and Professional Conduct'', 2018.
  \item (ISC)\textsuperscript{2}, ``(ISC)\textsuperscript{2} Code of Ethics''.
\end{enumerate}
}
